%-------------------------------------------------
% FileName: prac-sum.tex
% Author: Safin (zhaoqid@zsc.edu.cn)
% Version: 0.1
% Date: 2021-11-21
% Description: 实践总结
% Others: 
% History: origin
%------------------------------------------------- 

% 断页
\clearpage 

\chapter{实践总结}

% 本章，简单总结实践完成之后的想法
%实践总结主要按照下面的几部分内容展开，也可以加入自己的其它想法(成文是注释掉这段提示)。
%\section{实践之前对Linux编程的理解}
%\section{实践之后对Linux编程的理解}
%\section{对今后编程的想法}
%\section{对今后使用Linux系统的想法}
%\section{其它想法(如果没有就注释掉)} 
初学 Linux 编程时，我对它的认知极为狭隘，单纯将其等同于 “在命令行里敲代码、编译、运行” 的机械流程。
那时总觉得这种编程方式既繁琐又缺乏直观价值，完全体会不到背后的深层意义。尤其对于 FrameBuffer 这类底层接口，我仅停留在课本上 “抽象的图形驱动概念”，
从未想过它会与实际开发产生关联 —— 在我的固有认知里，绘图本就该调用 SDL、Qt 这类现成图形库，只需传入参数就能实现效果，
根本不需要接触系统底层的复杂知识，这种片面理解让我对 Linux 编程的核心价值视而不见。

直到参与这次实践课程，才彻底打破了原有认知。实践中，FrameBuffer 不再是抽象概念，而是需要亲手操作的 “显存入口”：为了实现雪花绘制，
我必须手动映射显存地址，将帧缓冲设备文件通过 mmap 映射到用户空间，直接对内存地址进行读写；根据帧缓冲的位深计算每个像素的存储方式，
手动拼接 RGB 颜色值并写入对应内存位置 —— 整个过程没有现成库的 “一键封装”，完全是 “从硬件层构建功能” 的硬核实践。这让我深刻意识到，Linux 编程的核心是 “与系统资源直接交互”：
它跳过了层层封装的黑盒，让开发者直面 CPU、内存、硬件设备等核心资源，让开发者看透技术的本质。

同时，这次实践也打破了 “图形开发必须依赖桌面框架” 的刻板印象。此前我一直认为，没有 GUI 环境就无法进行图形渲染，
而 FrameBuffer 的使用让我看到了另一种可能：在 Linux 系统中，仅凭操作显存就能实现雪花的分形绘制、动态下落等效果。
通过循环更新雪花坐标、逐像素刷新帧缓冲，最终呈现出流畅的动画，这让我理解了 “无 GUI 环境下的图形渲染逻辑”—— 本质上是通过直接操作硬件资源，
将像素数据按预期格式输出到显示设备，而桌面框架只是在此基础上做了更高层次的封装。这种认知的转变，让我对 “图形开发” 的理解不再局限于可视化工具，而是延伸到了系统底层的实现逻辑。

此次实践也为我后续的学习指明了方向：我不再满足于 “会用工具”，而是更注重 “底层原理 + 工具封装” 的结合。
比如这次手动实现 FrameBuffer 像素操作后，我计划将常用的绘图功能（如点、线、多边形绘制）封装成简易绘图函数库 —— 既保留对底层实现的清晰认知，
又能通过封装提升代码复用性，提高后续开发效率。

更重要的是，我逐渐明白 Linux 编程早已渗透到各类开发场景中，成为不可或缺的技能。以最近在嵌入式课程使用的 LVGL 图形库为例，
它不仅支持在 Linux 系统中基于 FrameBuffer 实现嵌入式 GUI 开发，并且在公司实习期间，项目开发时，我还知道了还能通过 Linux 环境搭建单片机的交叉编译工具链 —— 通过挂载操作系统镜像，
在 Linux 中编写、调试单片机程序，再通过串口将程序烧录到硬件中，实现对单片机内部逻辑的灵活修改。
无论是嵌入式设备的底层驱动开发、无 GUI 环境的图形应用开发，还是跨平台的工具链搭建，Linux 都扮演着核心角色。
未来，随着嵌入式、物联网等领域的发展，我接触的开发场景会更复杂，而 Linux 编程作为连接软件与硬件、跨越不同设备平台的桥梁，必然会成为我编程路上的 “必备技能”。

这次实践不仅让我掌握了 FrameBuffer 编程，更重要的是重塑了我对 Linux 编程的认知，
后续在实践中不断探索让 Linux 编程成为支撑我应对各类复杂开发场景的核心能力。
 


 